% Generated draft: requires mathematical and source review.
\documentclass[11pt]{article}
\usepackage[margin=25mm]{geometry}
\usepackage{fontspec}
\setmainfont{Latin Modern Roman}
\usepackage{amsmath,amssymb,mathtools,unicode-math}
\setmathfont{Latin Modern Math}
\usepackage{enumitem,needspace,xurl,hyperref}
\setlength{\parindent}{0pt}
\setlength{\parskip}{5pt}
\setlength{\emergencystretch}{3em}
\newcommand{\N}{\mathbb{N}}
\newcommand{\Z}{\mathbb{Z}}
\newcommand{\Q}{\mathbb{Q}}
\newcommand{\R}{\mathbb{R}}
\newcommand{\C}{\mathbb{C}}
\newcommand{\F}{\mathbb{F}}
\newcommand{\A}{\mathbb{A}}
\newcommand{\PP}{\mathbb{P}}
\newcommand{\E}{\mathbb{E}}
\newcommand{\eps}{\varepsilon}
\newcommand{\dd}{\,\mathrm{d}}
\begin{document}
\section*{UnsolvedMath ID: 30003009}
\textbf{Source problem code:} catalogue entry
\subsection{Definitions and mathematical statement}
Let \(n,k\in\Z\) satisfy \(n\geq3\) and \(2\leq k\leq n-1\). Under the standard convention, a convex body in \(\R^n\) is a compact convex set with nonempty interior. Let \(G(n,k)\) denote the Grassmannian of \(k\)-dimensional linear subspaces of \(\R^n\). For \(H\in G(n,k)\), let \(P_H:\R^n\to H\) be orthogonal projection, and let \(O(H)\) be the group of linear orthogonal transformations of \(H\).

\textbf{Reconstruction question, with an unresolved congruence qualification.} In the usual Euclidean meaning of congruent projections, the question is whether, for every pair of convex bodies \(K,L\subseteq\R^n\),
\[
\left[\forall H\in G(n,k)\ \exists U_H\in O(H)\ \exists b_H\in H:\quad P_HK=b_H+U_H(P_HL)\right]
\quad\Longrightarrow\quad
\left[\exists t\in\R^n\ \exists\varepsilon\in\{-1,1\}:\quad K=t+\varepsilon L\right].
\]
Here \(\varepsilon L=\{\varepsilon x:x\in L\}\), so the conclusion permits a translation of either \(L\) or its reflection through the origin. The projection congruence may depend on \(H\); the translation \(t\) and sign \(\varepsilon\) in the conclusion are global.

The source further requires that the congruences be ``induced by elements of \(O(k,H)\)'', but does not define this notation or specify whether translations within \(H\) are allowed. If it means solely linear orthogonal transformations of \(H\), the displayed antecedent must instead be
\[
\forall H\in G(n,k)\ \exists U_H\in O(H):\qquad P_HK=U_H(P_HL).
\]
These antecedents are distinct, and the supplied record does not determine which is intended. Both allow orientation-reversing orthogonal maps; no restriction to direct congruence is stated. Orthogonal projection, linear Grassmannians, and the full-dimensional convex-body convention above are standard interpretations rather than definitions explicitly supplied by the record.
\subsection{Short English statement}
For each \(2\leq k\leq n-1\), the question asks whether two convex bodies whose corresponding \(k\)-dimensional orthogonal projections are congruent must differ by a translation and possibly reflection through the origin. The undefined notation \(O(k,H)\) leaves unresolved whether the projection congruences may include translations. The supplied status is partially solved; this is a statement-only draft requiring mathematical review.
\subsection{Sources}
\begin{itemize}
\item Supplied local UnsolvedMath registry record, permanent ID \(30003009\), source problem code \texttt{catalogue entry}, title: \emph{Reconstruction of Convex Bodies from Congruent Projections}. The record gives proposed year 2015 and names no individual proposer.
\item Supplied report citation: \emph{Convex Geometry and its Applications (2016)}, Oberwolfach Reports, pp.~3179--3269, DOI \texttt{10.4171/owr/2015/56}. \url{https://doi.org/10.4171/owr/2015/56}. The supplied background describes Oberwolfach Reports as documenting mathematical workshops and their open-problem discussions.
\item Supplied literature reference: Alfonseca, Cordier, and Ryabogin on directly congruent projections. \url{https://arxiv.org/abs/1412.2727}.
\item Supplied literature reference labelled ``published reconstruction results''. \url{https://doi.org/10.1007/s11856-016-1394-6}.
\item Literature assessment dated exactly \texttt{2026-08-21}: the supplied curation labels the problem partially solved and reports that the general congruent-projection reconstruction question remains unresolved, with affirmative results for symmetric bodies, polytopes, and several additional non-symmetric cases. This is supplied dated curation derived from the OWR record, not an independently verified assessment. The precise hypotheses of those special-case results are not supplied.
\item The supplied statement-status label is \texttt{corrected\_verified}. Its accompanying qualification says that only safe notation, typography, and TeX-delimiter cleanup was applied, preserving hypotheses, quantifiers, and subquestions. This label does not establish independent mathematical or source verification of this draft.
\end{itemize}
\end{document}
